\documentclass[10pt,a4paper]{amsart}
\usepackage{amsmath,amssymb,amsthm,graphicx}
% Retained minimal preamble for the Stacks Project chapter more-morphisms.tex (More on Morphisms).
% Verbatim excerpt of the chapter preamble: Xy-pic setup (native preamble lines 14-18) and the native
% theorem declarations with the mathematical macros (native preamble lines 154-217). Copied mechanically;
% nothing added, removed or redefined. The wrapper class and the amsmath/amssymb/amsthm/graphicx base
% packages are supplied by the assembler, so the class/externaldocument part of the native preamble is omitted.
% For commutative diagrams we use Xy-pic
\usepackage[all]{xy}

% We use 2cell for 2-commutative diagrams.
\xyoption{2cell}
% Theorem environments.
%
\theoremstyle{plain}
\newtheorem{theorem}[subsection]{Theorem}
\newtheorem{proposition}[subsection]{Proposition}
\newtheorem{lemma}[subsection]{Lemma}

\theoremstyle{definition}
\newtheorem{definition}[subsection]{Definition}
\newtheorem{example}[subsection]{Example}
\newtheorem{exercise}[subsection]{Exercise}
\newtheorem{situation}[subsection]{Situation}

\theoremstyle{remark}
\newtheorem{remark}[subsection]{Remark}
\newtheorem{remarks}[subsection]{Remarks}

\numberwithin{equation}{subsection}

% Macros
%
\def\lim{\mathop{\mathrm{lim}}\nolimits}
\def\colim{\mathop{\mathrm{colim}}\nolimits}
\def\Spec{\mathop{\mathrm{Spec}}}
\def\Hom{\mathop{\mathrm{Hom}}\nolimits}
\def\Ext{\mathop{\mathrm{Ext}}\nolimits}
\def\SheafHom{\mathop{\mathcal{H}\!\mathit{om}}\nolimits}
\def\SheafExt{\mathop{\mathcal{E}\!\mathit{xt}}\nolimits}
\def\Sch{\mathit{Sch}}
\def\Mor{\mathop{\mathrm{Mor}}\nolimits}
\def\Ob{\mathop{\mathrm{Ob}}\nolimits}
\def\Sh{\mathop{\mathit{Sh}}\nolimits}
\def\NL{\mathop{N\!L}\nolimits}
\def\CH{\mathop{\mathrm{CH}}\nolimits}
\def\proetale{{pro\text{-}\acute{e}tale}}
\def\etale{{\acute{e}tale}}
\def\QCoh{\mathit{QCoh}}
\def\Ker{\mathop{\mathrm{Ker}}}
\def\Im{\mathop{\mathrm{Im}}}
\def\Coker{\mathop{\mathrm{Coker}}}
\def\Coim{\mathop{\mathrm{Coim}}}

% Boxtimes
%
\DeclareMathSymbol{\boxtimes}{\mathbin}{AMSa}{"02}

%
% Macros for moduli stacks/spaces
%
\def\QCohstack{\mathcal{QC}\!\mathit{oh}}
\def\Cohstack{\mathcal{C}\!\mathit{oh}}
\def\Spacesstack{\mathcal{S}\!\mathit{paces}}
\def\Quotfunctor{\mathrm{Quot}}
\def\Hilbfunctor{\mathrm{Hilb}}
\def\Curvesstack{\mathcal{C}\!\mathit{urves}}
\def\Polarizedstack{\mathcal{P}\!\mathit{olarized}}
\def\Complexesstack{\mathcal{C}\!\mathit{omplexes}}
% \Pic is the operator that assigns to X its picard group, usage \Pic(X)
% \Picardstack_{X/B} denotes the Picard stack of X over B
% \Picardfunctor_{X/B} denotes the Picard functor of X over B
\def\Pic{\mathop{\mathrm{Pic}}\nolimits}
\def\Picardstack{\mathcal{P}\!\mathit{ic}}
\def\Picardfunctor{\mathrm{Pic}}
\def\Deformationcategory{\mathcal{D}\!\mathit{ef}}
\usepackage{cleveref}
\crefname{originallemma}{original lemma}{original lemmas}
\Crefname{originallemma}{Original lemma}{Original lemmas}
\makeatletter
\newlabel{lemma-weighting-base-change}{{0F3B}{source}{}{sourcepointer-1}{}}
\newlabel{lemma-weighting-base-change@cref}{{[originallemma][0][]0F3B}{[1][0][]source}{}{}{}}
\makeatother
\crefname{originalsection}{original section}{original sections}
\Crefname{originalsection}{Original section}{Original sections}
\makeatletter
\newlabel{section-etale-localization}{{04HF}{source}{}{sourcepointer-2}{}}
\newlabel{section-etale-localization@cref}{{[originalsection][0][]04HF}{[1][0][]source}{}{}{}}
\makeatother
\title{A morphism carrying a strictly positive weighting is universally open}
\author{The Stacks Project Authors}
\date{stacks-project:more-morphisms.tex:local-snapshot-20261011; GNU FDL 1.2 or later}
\begin{document}
\maketitle
\noindent\textit{Editorial note.} Counted span is the author's own proof environment, source lines 22318-22334, immediately following the statement at 22311-22316.\par\medskip
\noindent\textit{Editorial note.} The proof cites one chapter lemma and one chapter section by \textbackslash{}ref; both are navigation pointers into the same chapter.\par\medskip
\noindent\textit{Editorial note.} The proof is short but every reduction is displayed; no step is deferred to the author's reader.\par\medskip
\noindent\textit{Allowed standard prerequisites (editorial).}\begin{itemize}
\item Invariance of weightings under base change. chapter lemma lemma-weighting-base-change, cited at proof line 22320. Reduces universal openness to openness of the morphism itself.
\item Etale localization of quasi-finite morphisms. chapter Section section-etale-localization, cited at proof line 22327. After passing to an etale neighbourhood of a point of the target, makes the morphism finite over an open neighbourhood of a chosen point.
\item Definition of a weighting and of its integral. chapter definition definition-weighting, used at proof line 22331. Supplies local constancy of the integral along the finite piece, which is what converts positivity of the weighting into openness of the image.
\end{itemize}\medskip
\ifdefined\hypertarget\hypertarget{sourcepointer-1}{}\fi
\noindent\textit{Source reference.} Original lemma 0F3B: sources/stacks-more-morphisms/pinned-original-source.tex, source lines 22173--22173. Navigation pointer only: the counted proof cites this statement by \textbackslash{}ref as a standing input; the deduction the counted proof has to make is carried out inside the counted span, and no part of that deduction is delegated to this pointer.\par\medskip
\ifdefined\hypertarget\hypertarget{sourcepointer-2}{}\fi
\noindent\textit{Source reference.} Original section 04HF: sources/stacks-more-morphisms/pinned-original-source.tex, source lines 11898--11898. Navigation pointer only: the counted proof cites this statement by \textbackslash{}ref as a standing input; the deduction the counted proof has to make is carried out inside the counted span, and no part of that deduction is delegated to this pointer.\par\medskip
\begin{lemma}
\label{lemma-weighting-universally-open}
Let $f : X \to Y$ be a locally quasi-finite morphism.
Let $w : X \to \mathbf{Z}$ be a weighting. If $w(x) > 0$
for all $x \in X$, then $f$ is universally open.
\end{lemma}

\begin{proof}
Since the property is preserved by base change, see
Lemma \ref{lemma-weighting-base-change}, it suffices
to prove that $f$ is open. Since we may also replace
$X$ by any open of $X$, it suffices to prove that $f(X)$
is open. Let $y \in f(X)$. Choose $x \in X$ with $f(x) = y$.
It suffices to prove that $f(X)$ contains an open neighbourhood
of $y$ and it suffices to do so after replacing $Y$ by an
\'etale neighbourhood of $y$. By \'etale localization of
quasi-finite morphisms, see Section \ref{section-etale-localization},
we may assume
there is an open neighbourhood $U \subset X$ of $x$
such that $\pi = f|_U : U \to Y$ is finite. Then
$\int_\pi w|_U$ is locally constant and has positive value at $y$.
Hence $\pi(U)$ contains an open neighbourhood of $y$
and the proof is complete.
\end{proof}

\end{document}
